% Part: proof-theory
% Chapter: proof-search
% Section: tableaux

\documentclass[../../../include/open-logic-section]{subfiles}

\begin{document}

\olfileid{pt}{ps}{tab}

\olsection{Tableau}

Sistem tableau yaiku sistem !!{proof} kang raket
sesambungané karo kalkulus sekuèn lan mligi cocog
kanggo panelusuran bukti, kanthi tangan utawa piranti
lunak. Ing tableau, !!a{proof} awujud wit !!{formula}s,
dudu wit sekuèn. Nanging, béda karo deduksi alamiah,
aturan tableau ngibarataké aturan kalkulus sekuèn.
!!{proof} tableau sejatiné panelusuran modhèl kang
cetha katon. Mula tableau sok diarani tableau
\emph{semantis} utawa \emph{wit kabeneran}.

Tableau \emph{mawa tandha} iku wit !!{formula}s mawa
tandha, kang bisa awujud $\sFmla{\True}{!A}$ utawa
$\sFmla{\False}{!A}$. Wit tuwuh mudhun. Ing wiwitané
ana !!{formula}s kang nemtokaké apa kang arep
!!{prove}. Tuladhané, kanggo !!{prove} sekuèn
$!A, !B \Sequent !C, !D$, tableau diwiwiti saka
$\sFmla{\True}{!A}, \sFmla{\True}{!B},
\sFmla{\False}{!C}, \sFmla{\False}{!D}$.
!!^a{structure} kang ndadèkaké $!A$ lan $!B$
bener, déning $!C$ lan $!D$ salah, yaiku
!!a{structure} kang nuduhaké yèn
$!A, !B \Sequent !C, !D$ ora valid.

Simpul godhong (paling ngisor) ing tableau bisa
ditambahi nganggo aturan pamranjangan cabang, kang
nambah simpul ing cabang kang padha utawa nambah loro
simpul sadulur ing sangisoré, saéngga cabang pecah
dadi loro cabang anyar. Aturan bisa ditrapaké marang
saben !!{formula} mawa tandha ing sadhuwuré simpul
iku ing cabang kang padha. Aturan pamranjangan cabang
kanggo sistem tableau klasik \Log{Tc} kapacak ing
\olref{tab:Tc}.

\subfile{rules-Tc}

Kaya kang katon, aturan-aturan mau padha karo aturan
\Log{G3c} kang diwalik arahé. Tandha $\True$
lan $\False$ nuduhaké yèn !!a{formula} dadi
rumus utama ing sisih kiwa utawa tengen
sekuèn, kanthi urutan mau.

Cabang tableau diarani \emph{ketutup} yèn ngemot
$\sFmla{\True}{!A}$ lan $\sFmla{\False}{!A}$
bebarengan, utawa yèn ngemot $\sFmla{\True}{\lfalse}$.
Yèn saben cabang ketutup, tableau
sakabèhé diarani ketutup. Tuladhané, iki tableau
ketutup kanggo $!A \lor (!B \land !C) \Sequent
(!A \lor !B) \land (!A \lor !C)$, kang diwiwiti
saka $\sFmla{\True}{!A \lor (!B \land !C)}$
lan $\sFmla{\False}{(!A \lor !B) \land (!A \lor !C)}$:
\begin{oltableau}
  [\sFmla{\True}{\formula{A} \lor (\formula{B} \land \formula{C})},
    just=\TAss, checked
    [\sFmla{\False}{(\formula{A} \lor \formula{B}) \land
        (\formula{A} \lor \formula{C})}, just=\TAss, checked
      [\sFmla{\True}{\formula{A}}, just={\TRule{\True}{\lor}[1]}
        [\sFmla{\False}{\formula{A} \lor \formula{B}},
          just={\TRule{\False}{\land}[2]}, checked
          [\sFmla{\False}{\formula{A}},
            just={\TRule{\False}{\lor}[4]}
            [\sFmla{\False}{\formula{B}},
              just={\TRule{\False}{\lor}[4]}, close]
          ]
        ]
        [\sFmla{\False}{\formula{A} \lor \formula{C}},
          just={\TRule{\False}{\land}[2]}, checked
          [\sFmla{\False}{\formula{A}},
            just={\TRule{\False}{\lor}[4]}
            [\sFmla{\False}{\formula{C}},
              just={\TRule{\False}{\lor}[4]}, close]
          ]
        ]
      ]
      [\sFmla{\True}{\formula{B} \land \formula{C}},
        just={\TRule{\True}{\lor}[1]}, checked
        [\sFmla{\False}{\formula{A} \lor \formula{B}},
          just={\TRule{\False}{\land}[2]}, checked
          [\sFmla{\True}{\formula{B}},
            just={\TRule{\True}{\land}[3]}, move by=2
            [\sFmla{\True}{\formula{C}},
              just={\TRule{\True}{\land}[3]}
              [\sFmla{\False}{\formula{A}},
                just={\TRule{\False}{\lor}[4]}
                [\sFmla{\False}{\formula{B}},
                  just={\TRule{\False}{\lor}[4]},close
                ]
              ]
            ]
          ]
        ]
        [\sFmla{\False}{\formula{A} \lor \formula{C}},
          just={\TRule{\False}{\land}[2]}, checked
          [\sFmla{\True}{\formula{B}},
            just={\TRule{\True}{\land}[3]}, move by=2
              [\sFmla{\True}{\formula{C}},
                just={\TRule{\True}{\land}[3]}
                [\sFmla{\False}{\formula{A}},
                  just={\TRule{\False}{\lor}[4]}
                  [\sFmla{\False}{\formula{C}},
                  just={\TRule{\False}{\lor}[4]},close
                ]
              ]
            ]
          ]
        ]
      ]
    ]
  ]
\end{oltableau}

Tandha cèntang nuduhaké yèn aturan kang cocog wis
ditrapaké marang !!a{formula} ing kabèh cabang
sangisoré. Lambang $\otimes$ nuduhaké cabang ketutup.
Nomer larik ing sawisé aturan nuduhaké !!{formula}s
mawa tandha kang ditrapi aturan iku (yaiku
!!{formula} ing cabang kang padha lan ing larik
kang dituduhaké).

Panelusuran bukti gampang ing tableau. Kita ngasilaké
tableau tataran demi tataran. Ing tataran~$0$,
tulis !!{formula}s wiwitan ing pérangan paling dhuwur.
Ing tataran~$n+1$, tumrap saben simpul godhong,
golekana rumus mawa tandha kang durung dicèntang
lan paling dhuwur. Trapna aturan pamranjangan cabang
marang rumus iku. Sawisé aturan ditrapaké ing saben
cabang kang ngemot rumus mawa tandha mau, cèntang
rumus iku. (Tuladha ing ndhuwur diasilaké nganggo
algoritma iki.)

!!{formula}s mawa tandha kang awujud
$\sFmla{\True}{\lforall}$ lan
$\sFmla{\False}{\lexists}$ ora tau dicèntang, mula
kudu ditangani kanthi cara mligi. Yèn ana, kita
kudu njamin yèn aturan $\TRule{\True}{\forall}$
lan $\TRule{\False}{\lexists}$ pungkasane ditrapaké
kanggo kabèh tèrma kang cocog. Iki bisa dijamin,
tuladhané, kanthi ngetrapaké aturan mau ing saben
tataran marang kabèh rumus kaya mangkono ing saben
cabang (sawisé ngurusi !!{formula} mawa tandha
paling dhuwur kang dudu wujud iki). Tèrma~$t$
kang dienggo yaiku tèrma kapisan~$t$ kang bisa
dibangun saka !!{constant}s kang wis katon ing
cabang (utawa $\Obj c_0$ yèn ora ana !!{constant})
lan !!{function}s kang katon ing !!{formula}s
wiwitan, lan kang !!{formula} mawa tandha cocog,
$\sFmla{\True}{!B(t)}$ utawa~$\sFmla{\False}{!B(t)}$,
durung katon ing cabang. Yèn ora ana tèrma anyar
kang cocog, aturan iki ora nambah simpul ing tataran
iku; aturan dicoba manèh nalika ana tèrma anyar.

Yèn panelusuran iki ora ngasilaké tableau ketutup,
ana cabang kabuka winates kang kabèh !!{formula}s
nonatomik kang bisa dicèntang wis dicèntang lan
kabèh instansi tèrma kang cocog kanggo rumus
kuantor kang ora dicèntang wis ana ing cabang,
utawa panelusuran ngasilaké
wit tanpa wates kanthi cabang winates ing saben
simpul, kang mesthi nduwèni cabang tanpa wates.
Kaya ing \olref[cpl]{sec}, kita bisa netepaké
!!a{structure}~$\Struct{M}$ saéngga yèn
$\sFmla{\True}{!A}$ katon ing cabang, mula
$\Sat{M}{!A}$, lan yèn $\sFmla{\False}{!A}$
katon, mula $\Sat/{M}{!A}$. (Jupuk $\Theta =
\Setabs{!A}{\sFmla{\True}{!A} \text{ ing cabang}}$
lan $\Xi =
\Setabs{!A}{\sFmla{\False}{!A} \text{ ing cabang}}$.)

Tableau kang diasilaké déning cara panelusuran iki
gampang diterjemahaké dadi bukti \Log{G3c} kanthi
mènèhi label sekuèn marang saben simpul. Yèn
rumus-rumus wiwitan cocog karo sekuèn
$\Gamma \Sequent \Delta$, wènèhana label
$\Gamma \Sequent \Delta$ marang saben rumus wiwitan.
Yèn cabang dipanjangaké kanthi ngetrapaké aturan
pamranjangan cabang marang !!{formula} mawa tandha
$\sFmla{\True}{!A}$, simpul godhongé diwènèhi label
$!A, \Pi \Sequent \Lambda$; yèn aturan ditrapaké
marang !!{formula} mawa tandha $\sFmla{\False}{!A}$,
simpul godhongé diwènèhi label
$\Pi \Sequent \Lambda, !A$. Ing papan kang tableau
ngetrapaké aturan \TRule{\True}{\star}, trapna
aturan \LeftR{\star} marang sekuèn (kanthi $!A$
minangka rumus utama); ing papan kang ngetrapaké
aturan \TRule{\False}{\star}, trapna aturan
\RightR{\star}. Premis aturan iku sekuèn-sekuèn
kang dadi label simpul cocog kang ditambahaké ing
tableau. Nalika tableau ketutup amarga ngemot
$\sFmla{\True}{!A}$ lan $\sFmla{\False}{!A}$,
label simpul godhong pungkasan awujud sekuèn
$!A, \Pi \Sequent \Lambda, !A$. Sawetara simpul
urut-urutan ing tableau bisa diwènèhi sekuèn kang
padha; busak duplikaté. Tuladhané, tableau ing
ndhuwur diterjemahaké dadi !!{proof} iki:
\[\small
\Axiom$!A \fCenter !A, !B$
\RightLabel{\RightR{\lor}}
\UnaryInf$!A \fCenter !A \lor !B$
\Axiom$!A \fCenter !A, !C$
\RightLabel{\RightR{\lor}}
\UnaryInf$!A \fCenter !A \lor !C$
\RightLabel{\RightR{\land}}
\BinaryInf$!A \fCenter (!A \lor !B) \land (!A \lor !C)$
\Axiom$!B, !C \fCenter !A, !B$
\RightLabel{\RightR{\lor}}
\UnaryInf$!B, !C \fCenter !A \lor !B$
\RightLabel{\LeftR{\land}}
\UnaryInf$!B \land !C \fCenter !A \lor !B$
\Axiom$!B, !C \fCenter !A, !C$
\RightLabel{\RightR{\lor}}
\UnaryInf$!B, !C \fCenter !A \lor !C$
\RightLabel{\LeftR{\land}}
\UnaryInf$!B \land !C \fCenter !A \lor !C$
\RightLabel{\RightR{\land}}
\BinaryInf$!B \land !C \fCenter (!A \lor !B) \land
(!A \lor !C)$
\RightLabel{\LeftR{\lor}}
\BinaryInf$!A \lor (!B \land !C) \fCenter (!A \lor !B) \land
(!A \lor !C)$
\DisplayProof
\]

\end{document}
